\section*{Capítulo \ref{ch:b2:meiosis-heredity} --- Meiosis, mezcla genética y herencia}

\begin{solution}{exo:b2:meiosis-heredity:1}
Mitosis: una división, sin apareamiento de homólogos, dos células, cada
una genéticamente idéntica a la madre ($2n \to 2n$). \omterm{def:b2:meiosis-heredity:meiosis}{Meiosis}: dos
divisiones tras una sola replicación, los homólogos se aparean y se
entrecruzan, cuatro células, \omterm{def:b2:meiosis-heredity:meiosis}{haploides} y todas genéticamente distintas
($2n \to n$).
\end{solution}

\begin{solution}{exo:b2:meiosis-heredity:2}
$2^{23} = 8.4\times 10^{6}$; $2^{4} = 16$; $2^{7} = 128$ --- antes de
que el entrecruzamiento multiplique cada una de estas cifras sin límite.
\end{solution}

\begin{solution}{exo:b2:meiosis-heredity:3}
$Tt\times Tt$: \omterm{def:b2:meiosis-heredity:vocabulary}{genotipos} $1\,TT : 2\,Tt : 1\,tt$, \omterm{def:b2:meiosis-heredity:vocabulary}{fenotipos} $3$ altas
: $1$ enana. $Tt\times tt$: $1\,Tt : 1\,tt$, $1 : 1$. Se hace un
cruzamiento de prueba de la planta alta con una enana: si todas son
altas, es $TT$; si la mitad son enanas, es $Tt$ (o se autofecunda: $TT$
se reproduce sin variación y $Tt$ segrega $3 : 1$).
\end{solution}

\begin{solution}{exo:b2:meiosis-heredity:4}
Los genes ligados están en el mismo cromosoma y sus combinaciones
parentales de \omterm{def:b2:meiosis-heredity:vocabulary}{alelos} están sobrerrepresentadas entre los gametos; la
\omterm{def:b2:meiosis-heredity:linkage}{frecuencia de recombinación} es la fracción de gametos recombinantes
(descendientes recombinantes de un cruzamiento de prueba); un \omterm{def:b2:meiosis-heredity:linkage}{centimorgan}
es un uno por ciento de recombinación. Los genes muy alejados en un
cromosoma largo tienen al menos un entrecruzamiento entre ellos en casi
todas las \omterm{def:b2:meiosis-heredity:meiosis}{meiosis}, y los entrecruzamientos múltiples aleatorizan las
combinaciones, de modo que $r$ alcanza $\tfrac12$ --- indistinguible de
la \omterm{thm:b2:meiosis-heredity:mixing}{segregación independiente}.
\end{solution}

\begin{solution}{exo:b2:meiosis-heredity:5}
Total 556; esperados $312.75 : 104.25 : 104.25 : 34.75$; $\chi^2 =
2.25^2/312.75 + 3.75^2/104.25 + 3.25^2/104.25 + 2.75^2/34.75 = 0.016
+ 0.135 + 0.101 + 0.218 = 0.47 < 7.81$: compatible con $9 : 3 : 3 :
1$.
\end{solution}

\begin{solution}{exo:b2:meiosis-heredity:6}
Frente a $1 : 1 : 1 : 1$ (250 de cada): $\chi^2 = (162^2 + 138^2 + 147^2
+ 153^2)/250 = 361 \gg 7.81$: ligados. Recombinantes: $103 + 97 = 200$
de 1000: $r = 0.20$, \qty{20}{cM}. En repulsión, $Ab/aB$ da 400
$Ab$, 400 $aB$, 100 $AB$ y 100 $ab$: la misma distancia, con las clases
parentales y recombinantes intercambiadas.
\end{solution}

\begin{solution}{exo:b2:meiosis-heredity:7}
$d = -\tfrac12\ln(1 - 2r)$: $r = 0.20$ da $-\tfrac12\ln 0.6 =
0.255$, \qty{25.5}{cM}; $r = 0.45$ da $-\tfrac12\ln 0.1 = 1.15$,
\qty{115}{cM}. $r = \tfrac12(1 - e^{-2d})$: \qty{30}{cM} dan
$\tfrac12(1 - e^{-0.6}) = 0.226$; \qty{100}{cM} dan $\tfrac12(1 -
e^{-2}) = 0.43$. Por debajo de \qty{10}{cM} la corrección es
despreciable; por encima de \qty{30}{cM} es imprescindible, y por encima
de \qty{50}{cM} $r$ se satura, de modo que un único cruzamiento de dos
puntos ya no puede medir la distancia --- los genes lejanos se
cartografían encadenando genes cercanos.
\end{solution}

\begin{solution}{exo:b2:meiosis-heredity:8}
(a) $X^wX^w \times X^+Y$: todos los hijos blancos ($X^wY$), todas las
hijas rojas ($X^+X^w$). (b) $X^+X^w \times X^wY$: hijos, la mitad rojos y
la mitad blancos; hijas, la mitad rojas ($X^+X^w$) y la mitad blancas
($X^wX^w$). (c) hijas
$X^+X^w$ $\times$ hermanos $X^wY$: como en (b) --- la mitad de cada sexo
blanca, la mitad roja.
\end{solution}

\begin{solution}{exo:b2:meiosis-heredity:9}
Dos enzimas en serie fabrican el pigmento; el púrpura necesita un \omterm{def:b2:meiosis-heredity:vocabulary}{alelo}
\omterm{def:b2:meiosis-heredity:vocabulary}{dominante} en ambos loci ($C\_P\_$). Las líneas son $CCpp$ y $ccPP$ (cada
una blanca por una razón distinta); el híbrido $CcPp$ es púrpura;
autofecundado da $9\,C\_P\_$ púrpuras por $7$ blancas ($3\,C\_pp + 3\,ccP\_ +
1\,ccpp$). Cruzamiento de prueba $CcPp\times ccpp$: $1$ púrpura
($CcPp$) : $3$
blancas.
\end{solution}

\begin{solution}{exo:b2:meiosis-heredity:10}
Parentales $+++$, $abc$; dobles entrecruzamientos $+b+$, $a+c$; $b$ es
el gen que ha cambiado: orden $a$--$b$--$c$. $r_{ab} = (11 + 9 + 2)/1202
= 0.018$ (\qty{1.8}{cM}); $r_{bc} = (118 + 112 + 2)/1202 = 0.193$
(\qty{19.3}{cM}). Dobles esperados $0.018\times 0.193\times 1202 =
4.2$; observados 2: coincidencia $0.47$, interferencia $0.53$.
\end{solution}

\begin{solution}{exo:b2:meiosis-heredity:11}
Ascos de segunda división: $300/1000 = \qty{30}{\%}$; distancia $=
\tfrac12\times 30 = \qty{15}{cM}$. Un entrecruzamiento entre el gen y el
centrómero implica solo a dos de las cuatro cromátidas, de modo que un
asco con segregación en la segunda división lleva dos cromátidas
recombinantes y dos parentales: la \omterm{def:b2:meiosis-heredity:linkage}{frecuencia de recombinación} es la
mitad de la frecuencia de esos ascos. Para $2 : 4 : 2$: un
entrecruzamiento entre el locus y el centrómero intercambia las dos
cromátidas internas, de modo que cada semihuso de la \omterm{def:b2:meiosis-heredity:meiosis}{meiosis} I lleva una
cromátida $A$ y una $a$; la \omterm{def:b2:meiosis-heredity:meiosis}{meiosis} II las separa entonces en el orden
$A\,a \mid
a\,A$ (o el inverso), y la mitosis final duplica cada espora:
$2 : 4 : 2$.
\end{solution}

\begin{solution}{exo:b2:meiosis-heredity:12}
Su madre es portadora obligada (recibió el único X de su padre), de modo
que la mujer recibió el \omterm{def:b2:meiosis-heredity:vocabulary}{alelo} con probabilidad $\tfrac12$. Dos hijos
varones sanos: una portadora tiene una probabilidad
$(\tfrac12)^2 = \tfrac14$ de tener dos hijos sanos; una no portadora,
una probabilidad de 1. Bayes:
$P(\text{portadora}\mid\text{datos}) = (\tfrac12\times\tfrac14)/(\tfrac12
\times\tfrac14 + \tfrac12\times 1) = \tfrac{1/8}{5/8} = \tfrac15$.
\end{solution}

\begin{solution}{pb:b2:meiosis-heredity:1}
\textbf{1.} Progenitores $vg/vg\;+/+$ y $+/+\;e/e$; $F_1$ $+/vg\;+/e$;
gametos $+\,+$, $+\,e$, $vg\,+$, $vg\,e$, cada uno $\tfrac14$.
\textbf{2.} $900 : 300 : 300 : 100$.
\textbf{3.} $\chi^2 = 8^2/900 + 9^2/300 + 4^2/300 + 3^2/100 = 0.07 +
0.27 + 0.05 + 0.09 = 0.48 < 7.81$: se cumple la \omterm{thm:b2:meiosis-heredity:mixing}{segregación
independiente}.
\textbf{4.} Las homocigotas silvestres en ambos loci son $\tfrac{1}{16}$
del total, y $\tfrac{9}{16}$ son silvestres: $\tfrac19$.
\textbf{5.} Silvestre, vestigial, ébano y vestigial ébano, $\tfrac14$
cada una.
\textbf{6.} El \omterm{def:b2:meiosis-heredity:vocabulary}{fenotipo} es producto del \omterm{def:b2:meiosis-heredity:vocabulary}{genotipo} y del ambiente (la
temperatura actúa sobre las enzimas del pigmento): un \omterm{def:b2:meiosis-heredity:vocabulary}{genotipo} especifica
una norma de reacción, no un carácter fijo.
\textbf{7.} $b\,+/+\,vg$ (repulsión): $b$ llegó con $+_{vg}$ de un
progenitor, y $vg$ con $+_b$ del otro.
\textbf{8.} Parentales: negras ($b\,+$) 965 y vestigiales ($+\,vg$) 944;
recombinantes: negras vestigiales ($b\,vg$) 206 y silvestres ($+\,+$) 185.
El cromosoma silvestre $+\,+$ no existía en la $F_1$; solo puede
formarse por entrecruzamiento.
\textbf{9.} $r = 391/2300 = 0.17$: \qty{17}{cM}.
\textbf{10.} $d = -\tfrac12\ln(1 - 0.34) = -\tfrac12\ln 0.66 = 0.21$:
\qty{21}{cM}.
\textbf{11.} En los machos de \emph{Drosophila} no hay entrecruzamiento:
los gametos del macho solo llevan los dos cromosomas parentales.
\textbf{12.} $F_1$ $b\,vg/+\,+$ (acoplamiento): clases parentales
silvestre y negra vestigial, un \qty{41.5}{\%} cada una; recombinantes
negra y vestigial, un \qty{8.5}{\%} cada una.
\textbf{13.} $X^hX^+$: una hija recibe el único X de su padre, que lleva
$h$, y un X normal de su madre no afectada.
\textbf{14.} $\tfrac12$ (varón) $\times$ $\tfrac12$ (recibe $X^h$) $=
\tfrac14$.
\textbf{15.} $\tfrac12$: recibe uno cualquiera de los dos X de la madre;
el padre aporta un X normal.
\textbf{16.} Cada hijo queda no afectado con probabilidad $\tfrac34$:
$(\tfrac34)^3 = \tfrac{27}{64} = 0.42$.
\textbf{17.} Un hijo varón recibe el Y de su padre, nunca su X. Una mujer
solo está afectada si es $X^hX^h$: un padre afectado y una madre
portadora (o afectada), lo cual es raro.
\textbf{18.} Probabilidad previa: $\tfrac12$. Dos hijos sanos tienen
probabilidad $\tfrac14$ si es portadora, y 1 si no lo es: probabilidad
posterior $(\tfrac12\times
\tfrac14)/(\tfrac12\times\tfrac14 + \tfrac12) = \tfrac15$.
\textbf{19.} El único X de un hijo varón procede de su madre, y el Y no
lleva ninguno de estos genes, de modo que cada hijo muestra sin
enmascarar un gameto materno: el cruzamiento es su propio cruzamiento de
prueba.
\textbf{20.} Parentales $+++$ (853), $ywm$ (833); dobles entrecruzamientos
$+w+$ (2), $y+m$ (2); ha cambiado $w$: orden $y$--$w$--$m$.
\textbf{21.} $r_{yw} = (24 + 26 + 2 + 2)/2000 = 0.027$; $r_{wm} = (128
+ 132 + 2 + 2)/2000 = 0.132$. Mapa: $y$ --- \qty{2.7}{cM} --- $w$ ---
\qty{13.2}{cM} --- $m$.
\textbf{22.} Esperados $0.027\times 0.132\times 2000 = 7.1$; observados
4: coincidencia $0.56$, interferencia $0.44$.
\textbf{23.} Recombinantes $y$--$m$: $y++$, $+wm$, $yw+$, $++m$: $(24 +
26 + 128 + 132)/2000 = 0.155$, frente a $0.027 + 0.132 = 0.159$: los
cuatro dobles entrecruzamientos restablecen la combinación parental
$y$--$m$ y pasan desapercibidos, $2\times 4/2000 = 0.004$.
\textbf{24.} Las hijas reciben el X $y\,w\,m$ del padre y un X materno:
todas son heterocigotas y de \omterm{def:b2:meiosis-heredity:vocabulary}{fenotipo} silvestre, lleven el X recombinante
que lleven; la recombinación es invisible.
\textbf{25.} $y$ --- \qty{2.7}{cM} --- $w$ --- \qty{13.2}{cM} --- $m$;
coeficiente de coincidencia $0.56$, interferencia $0.44$.
\end{solution}
